% STATUS: written 2026-11-06 (English gate clean; sources read in full as recorded below).
% Sources read in full for the cited passages: Ferrero-Lequain-Nowicki 1992 (pp. 45--49), Daigle-Freudenburg 1998 (pp. 353--359), van den Essen-Kuroda-Koenig 2021 (pp. 65--66), Venegas-Venegas 2025 (pp. 1--3, 7), Cid-Veloso 2026 (pp. 1--4).
% Sources: FerreroLequainNowicki1992, DaigleFreudenburg1998UFD, VanEssenKurodaKoenig2021,
% VenegasVenegas2025Cancellation, CidVeloso2026Isotropy, ChitayatDaigleRegeta2026,
% Tanaka2010ModuleDerivations, CidVeloso2026TameIsotropy, PetravchukShevchykSysak2017,
% Zaidenberg2026LocallyFinite.
\section{Locally nilpotent derivations}\label{sec:commutative:locally-nilpotent}

Section~\ref{sec:commutative:zariski-lipman} asked what the algebra feels when
its derivations form a free module. This section asks a sharper
question~\cite[p.~354]{DaigleFreudenburg1998UFD}: what happens when a
derivation behaves like a nilpotent operator on every element?

Chapter~2 defined the class~\cite[Section~1.1.9, p.~12]{Freudenburg2017}. A
derivation $D$ is \emph{locally nilpotent} when every element $b$ is killed by
some power of $D$ \cite[Section~1.1.9, p.~12]{Freudenburg2017}. For such a
derivation the map $\exp D$ is well
defined~\cite[Section~1.1.9, p.~12]{Freudenburg2017}.

On a polynomial ring these derivations have a geometric
meaning~\cite[p.~354]{DaigleFreudenburg1998UFD}. Algebraic actions of the
additive group $\mathbb{G}_a$ on affine space $\mathbb{A}^n$ correspond to
locally nilpotent derivations~\cite[p.~354]{DaigleFreudenburg1998UFD}.

So the question of the section is
this~\cite[Theorem, p.~47]{FerreroLequainNowicki1992}. When does a derivation
built from locally nilpotent pieces stay locally
nilpotent~\cite[Theorem, p.~47]{FerreroLequainNowicki1992}? And what do these
derivations look
like~\cite[Theorem, p.~47]{FerreroLequainNowicki1992}?

Ferrero, Lequain and Nowicki answer the first part for a combination of two
commuting derivations \cite[Theorem, p.~47]{FerreroLequainNowicki1992}.

\begin{theorem}[{\cite[Theorem, p.~47]{FerreroLequainNowicki1992}}]\label{thm:comm:fln}
Let $R$ be a commutative, reduced, $\mathbb{Z}$-torsion free ring, let $d$
and $\delta$ be two locally nilpotent derivations of $R$ which commute, and
let $b\in R$ be such that $\delta(b) = d(b) = 0$. Let
$\Delta = a\,d + b\,\delta$ with $a\in R$. Then $\Delta$ is locally nilpotent
if and only if $d(a) = 0$.
\end{theorem}

For the polynomial ring in two variables the criterion becomes a clean
description. It lists all locally nilpotent derivations of a triangular form
\cite[Corollary~6, p.~49]{FerreroLequainNowicki1992}.

\begin{corollary}[{\cite[Corollary~6, p.~49]{FerreroLequainNowicki1992}}]\label{cor:comm:fln6}
Let $K[x, y]$ be the polynomial ring in two indeterminates over a reduced,
$\mathbb{Z}$-torsion free ring $K$. Let $f \in K[x, y]$ and $b \in K$, and
let $\Delta$ be the $K$-derivation of $K[x, y]$ defined by $\Delta(x) = f$
and $\Delta(y) = b$. Then $\Delta$ is locally nilpotent if and only if
$f \in K[y]$.
\end{corollary}

In dimension two much more is true
\cite[pp.~356--357]{DaigleFreudenburg1998UFD}. The theorem of Rentschler from
1968 describes every locally nilpotent derivation of the plane: its kernel is
generated by one coordinate
\cite[pp.~356--357]{DaigleFreudenburg1998UFD}. Daigle and Freudenburg recall
it in Jacobian form, where for $P\in L[X,Y]$ they write $D_P$ for the
derivation $h \mapsto P_X h_Y - P_Y h_X$
\cite[pp.~356--357]{DaigleFreudenburg1998UFD}.

\begin{theorem}[Rentschler; {\cite[pp.~356--357]{DaigleFreudenburg1998UFD}}]\label{thm:comm:rentschler}
Let $L$ be a field of characteristic zero and let $D$ be a nonzero locally
nilpotent derivation of $L[X, Y]$. Then there exist $P, Q$ such that
$L[X, Y] = L[P, Q]$ and $\ker D = L[P]$. Moreover, there exists $a\in L[P]$
such that $D h = a\,(P_X h_Y - P_Y h_X)$ for all $h \in L[X, Y]$.
\end{theorem}

Then Daigle and Freudenburg extend the theorem of Rentschler. Their version
works over a factorial base ring in place of a field
\cite[Theorem~2.4, pp.~358--359]{DaigleFreudenburg1998UFD}.

\begin{theorem}[{\cite[Theorem~2.4, pp.~358--359]{DaigleFreudenburg1998UFD}}]\label{thm:comm:df24}
Let $R$ be a unique factorisation domain containing $\mathbb{Q}$, let
$B = R[X, Y]$ and let $K$ be the field of fractions of $R$. For a nonzero
$R$-derivation $D$ of $B$, the following are equivalent: $D$ is locally
nilpotent; and $D = a\,D_P$ for some $P \in B$ which is a variable of
$K[X, Y]$ with $\gcd_B(P_X, P_Y) = 1$ and some
$a \in R[P]\setminus\{0\}$. In this case $\ker D = R[P]$.
\end{theorem}

Beyond dimension two, coordinates no longer classify the
kernels~\cite[p.~65]{VanEssenKurodaKoenig2021}. One measures them instead, and
van den Essen, Kuroda and Koenig explain how the measure
arose~\cite[p.~65]{VanEssenKurodaKoenig2021}. Their work produced a list of
hypersurfaces that might be
counterexamples~\cite[p.~65]{VanEssenKurodaKoenig2021}.

The best known is a hypersurface $V$ in $\mathbb{C}^4$, given by
$x + x^2 y + z^2 + t^3 = 0$
\cite[p.~65]{VanEssenKurodaKoenig2021}. It is smooth, factorial and
contractible, and diffeomorphic to $\mathbb{R}^6$
\cite[p.~65]{VanEssenKurodaKoenig2021}. To separate such candidates from affine
space, one meets all locally nilpotent derivations in their
kernels~\cite[p.~65]{VanEssenKurodaKoenig2021}.

\begin{definition}[{\cite[Definition~3.1.1, p.~65]{VanEssenKurodaKoenig2021}}]\label{def:comm:ml}
Let $R$ be a ring of characteristic zero. The \emph{Makar-Limanov invariant}
of $R$ is the intersection of the kernels of all locally nilpotent
derivations of $R$:
\[
\operatorname{ML}(R) = \bigcap_{D \in \operatorname{LND}(R)} \ker D .
\]
If $W$ is an algebraic variety with coordinate ring $A(W)$, one sets
$\operatorname{ML}(W) = \operatorname{ML}(A(W))$.
\end{definition}

For a field $k$ of characteristic zero the invariant of the polynomial ring is
$k$ \cite[Example~3.1.2, p.~66]{VanEssenKurodaKoenig2021}, and that holds in
every number of variables.

Makar-Limanov showed that the element $x$ of the Koras-Russell hypersurface
lies in every kernel~\cite[p.~66]{VanEssenKurodaKoenig2021}. It turned out
that $\operatorname{ML}(V) = \mathbb{C}[x]$
\cite[p.~66]{VanEssenKurodaKoenig2021}. The invariant of affine three-space is
$\mathbb{C}$ \cite[p.~66]{VanEssenKurodaKoenig2021}. So the hypersurface $V$
is not isomorphic to $\mathbb{C}^3$ \cite[p.~66]{VanEssenKurodaKoenig2021}.

Chapter~12 returns to this invariant for polynomial
rings~\cite[p.~65]{VanEssenKurodaKoenig2021}. The cancellation problem behind
these invariants is treated in the cited
book~\cite[p.~65]{VanEssenKurodaKoenig2021}.

The invariant is the tool of the cancellation problem, and a recent survey
follows that use through three settings
\cite[pp.~1--3]{VenegasVenegas2025Cancellation}. The survey states the
problem of Zariski in the commutative case, then in noncommutative algebras,
then in skew extensions \cite[pp.~1--3]{VenegasVenegas2025Cancellation}. The
key technical input there is the stability of the invariant under adjoining a
variable. For a two-dimensional unique factorization domain it gives a
dichotomy \cite[Theorem~2.2, p.~7]{VenegasVenegas2025Cancellation}.

\begin{theorem}[{\cite[Theorem~2.2, p.~7]{VenegasVenegas2025Cancellation}}]\label{thm:comm:ml-dichotomy}
Let $k$ be an algebraically closed field of characteristic zero and let $R$
be a two-dimensional affine $k$-domain that is a unique factorization
domain. Then exactly one of the following holds: $\operatorname{ML}(R) = R$,
or $\operatorname{ML}(R) = k$ and $R \cong k[x,y]$.
\end{theorem}

So in dimension two the invariant separates the rigid algebras from the
polynomial ring itself. The survey reads the cancellativity of the affine
plane off this dichotomy \cite[Theorem~2.3, p.~7]{VenegasVenegas2025Cancellation}.
It also reports where the method fails: in the skew case the stability of the
invariant can break down
\cite[pp.~2--3]{VenegasVenegas2025Cancellation}.

Beyond cancellation, the locally nilpotent derivations of the plane have
constrained symmetry. Cid and Veloso compute the group of automorphisms that
commute with a locally finite derivation \cite[p.~1]{CidVeloso2026Isotropy}.
Every nonzero locally finite derivation of $K[X,Y]$ has a nontrivial such
group \cite[Proposition~3.1, p.~3]{CidVeloso2026Isotropy}.

\begin{theorem}[{\cite[Proposition~3.1, p.~3]{CidVeloso2026Isotropy}}]\label{thm:comm:isotropy-nontrivial}
Let $K$ be an algebraically closed field of characteristic zero and let $B$
be an affine $K$-algebra. If $D \in \operatorname{LFD}(B)$ with $D \neq 0$,
then the group of automorphisms of $B$ that commute with $D$ is nontrivial.
\end{theorem}

The proof takes the exponential $\exp(tD)$, a well defined automorphism
because $D$ is locally finite
\cite[p.~4]{CidVeloso2026Isotropy}. For each $t$ this exponential commutes
with $D$, and for some $t$ it differs from the identity
\cite[p.~4]{CidVeloso2026Isotropy}. So a simple derivation of $K[X,Y]$ cannot
be locally finite, because a simple derivation of the plane has a trivial
commuting group \cite[p.~4]{CidVeloso2026Isotropy}. Theorem~\ref{thm:pw:isotropy-simple}
of Chapter~12 states that triviality. The same authors compute the full
commuting group for each of the four normal forms of a locally finite
derivation of $K[X,Y]$ \cite[Theorems~3.3--3.9, pp.~4--8]{CidVeloso2026Isotropy}.

The same authors also compute the tame part of the group, generated by the
elementary automorphisms that commute with $D$
\cite[p.~1]{CidVeloso2026TameIsotropy}. For a locally nilpotent derivation the tame
group and the tame group of its exponential agree
\cite[Proposition~4.1, p.~9]{CidVeloso2026TameIsotropy}.

\begin{theorem}[{\cite[Proposition~4.1, p.~9]{CidVeloso2026TameIsotropy}}]\label{thm:comm:tame-isotropy-exp}
Let $D$ be a locally nilpotent derivation of $K[X,Y]$. Then the tame
isotropy group of $D$ equals the tame isotropy group of
$\exp(D)$.
\end{theorem}

For a derivation with a nontrivial semisimple part the two full groups can
differ, while the tame groups still agree
\cite[pp.~1--2]{CidVeloso2026TameIsotropy}.

Bianchi, Freitas and Veloso introduce the same tame group for triangular
derivations in up to three variables and for simple Shamsuddin derivations in
any number of variables \cite[Abstract, p.~1]{BianchiFreitasVeloso2025Tame}.

Chunaev and Evdokimova study the subgroup of locally nilpotent derivations
that commute with a homogeneous one. They find its maximal homogeneous
subgroups on affine trinomial hypersurfaces, and give when such a subgroup is
maximal \cite[Abstract, p.~1]{ChunaevEvdokimova2026}.

A derivation is locally finite when the vector space spanned by its powers on
each element is finite dimensional
\cite[Definition~2.1, p.~2]{ChitayatDaigleRegeta2026}. A set of derivations
is locally finite when one finite-dimensional space serves every member of
the set, and weakly locally finite when every finite subset is
\cite[Definition~2.1, p.~2]{ChitayatDaigleRegeta2026}. Chitayat, Daigle and
Regata ask when a Lie algebra of locally finite derivations is weakly
locally finite \cite[pp.~1--2]{ChitayatDaigleRegeta2026}. They answer for
the solvable case on a quasi-affine variety
\cite[Corollary~4, p.~2]{ChitayatDaigleRegeta2026}.

\begin{theorem}[{\cite[Corollary~4, p.~2]{ChitayatDaigleRegeta2026}}]\label{thm:comm:locally-finite-solvable}
Let $X$ be a quasi-affine variety over a field $k$. If $L$ is a solvable
Lie subalgebra of $\operatorname{Der}_k \mathcal{O}_X(X)$ whose elements
are locally finite derivations, then $L$ is weakly locally finite. Moreover,
$L$ admits a countable increasing filtration by locally finite Lie
subalgebras.
\end{theorem}

The proof reduces the question to a derivation-finite algebra, one fixed by
a finite set of values \cite[p.~2]{ChitayatDaigleRegeta2026}. The algebra of
regular functions on a quasi-affine variety is of this kind
\cite[p.~4]{ChitayatDaigleRegeta2026}. So the result applies to the
coordinate rings of the varieties met in this chapter.

The same local nilpotency carries over to modules, and there the finite
generation can fail. Tanaka works with a finitely generated module $M$ over
an affine algebra $B$, with a module derivation $\delta_{M}$ that is
locally nilpotent \cite[p.~1]{Tanaka2010ModuleDerivations}. Let $A$ be the
kernel of $\delta$ on $B$ and $M_{0}$ the kernel of $\delta_{M}$ on $M$
\cite[p.~1]{Tanaka2010ModuleDerivations}. He asks whether $M_{0}$ is a
finitely generated $A$-module \cite[p.~2]{Tanaka2010ModuleDerivations}.

\begin{theorem}[{\cite[Theorem~4.6, p.~10]{Tanaka2010ModuleDerivations}}]\label{thm:comm:module-kernel}
Let $\delta$ be a locally nilpotent derivation of an affine $k$-algebra $B$
and let $M$ be a finitely generated torsion-free $(B,\delta)$-module.
Suppose that $A=\ker\delta$ is a noetherian domain. Then $M_{0}=\ker\delta_{M}$
is a finitely generated $A$-module.
\end{theorem}

The theorem settles the question in dimension at most three, since the
kernel of a locally nilpotent derivation in that range is finitely generated
\cite[Corollary~4.7, p.~10]{Tanaka2010ModuleDerivations}. In higher
dimension the answer can be negative, and the counterexamples come from the
counterexamples to the fourteenth problem of Hilbert
\cite[pp.~2--3]{Tanaka2010ModuleDerivations}. So the module question is as
hard as the invariant question behind it.

Locally nilpotent derivations also form Lie algebras of their own, and
Petravchuk, Shevchyk and Sysak study the locally nilpotent subalgebras of
the derivations of an integral domain \cite[p.~1]{PetravchukShevchykSysak2017}.
They measure such a subalgebra by its rank over the fraction field
\cite[p.~1]{PetravchukShevchykSysak2017}.

\begin{theorem}[{\cite[Theorem~1, p.~5]{PetravchukShevchykSysak2017}}]\label{thm:comm:locally-nilpotent-series}
Let $A$ be an integral domain over a field $K$ of characteristic zero, let
$R$ be its fraction field, and let $L$ be a locally nilpotent subalgebra of
$R\operatorname{Der}_{K}A$ of rank $n$ over $R$. Then $L$ contains a series
of ideals $0=L_{0}\subset L_{1}\subset\cdots\subset L_{n}=L$ with
$\operatorname{rk}_{R}L_{i}=i$ for all $i$, and every quotient
$L_{i}/L_{i-1}$ is abelian.
\end{theorem}

So a locally nilpotent Lie algebra of derivations is built from abelian
pieces, one for each unit of rank. The result holds for the triangular
derivations of a polynomial ring, which are the model case
\cite[p.~2]{PetravchukShevchykSysak2017}.

On the affine plane the local finiteness can be read off the algebra
itself. Zaidenberg takes a solvable Lie subalgebra of the Lie algebra of
the automorphism group of the plane, generated by locally finite
derivations \cite[p.~1]{Zaidenberg2026LocallyFinite}. He proves that it is
triangulable and weakly locally finite
\cite[Theorem~1.2, p.~2]{Zaidenberg2026LocallyFinite}.

\begin{theorem}[{\cite[Theorem~1.2, p.~2]{Zaidenberg2026LocallyFinite}}]\label{thm:comm:planar-solvable}
Let $\mathfrak{h}=\langle\partial_{1},\dots,\partial_{k},\dots\rangle_{\mathrm{Lie}}$
be a solvable Lie subalgebra of the Lie algebra of the automorphism group
of the affine plane $\mathbb{A}^{2}$, generated by locally finite
derivations $\partial_{i}$. Then $\mathfrak{h}$ is triangulable and weakly
locally finite. If $\mathfrak{h}$ is generated by a finite set of locally
finite derivations, then it is locally finite.
\end{theorem}

The proof uses Newton polygons on the plane
\cite[pp.~10--16]{Zaidenberg2026LocallyFinite}. So the plane answers the
question that Corollary~\ref{thm:comm:locally-finite-solvable} answers in
general, and it adds triangulability, which the general statement does not
give.

Locally nilpotent derivations of the polynomial ring also leave two questions
open, and this section states both. % claim-ok: own signpost introducing the two moved problems, each carrying its own citation

\begin{problem}[{\cite[Section~3.9, p.~108]{Freudenburg2017}}]\label{prob:open:hdp}
Let $D$ be a standard homogeneous locally nilpotent derivation of
$k[x_1,\dots,x_n]$ with $D^2x_i=0$ for every $i$.
Is the rank of $D$ always strictly less than $n$?
\end{problem}

Freudenburg reports that this problem is open only in dimension $5$
\cite[p.~108]{Freudenburg2017}. A companion question asks whether a
homogeneous locally nilpotent derivation of degree two can have full rank, and
that question is open in every dimension \cite[Remark~3.51, p.~112]{Freudenburg2017}.

\begin{problem}[{\cite[Question~11.4, p.~288]{Freudenburg2017}}]\label{prob:open:bass}
Let a unipotent group $G$ act on the affine space $\mathbb{A}^n_k$.
Can the action be triangularized in the rational sense, that is, can one write
$k\langle x_1,\dots,x_n\rangle=k\langle y_1,\dots,y_n\rangle$ so that each
subfield $k\langle y_1,\dots,y_i\rangle$ is invariant under $G$?
\end{problem}

Bass asked this question in 1984, and Freudenburg records that it is still
open in general \cite[p.~288]{Freudenburg2017}. The same epilogue asks whether
the kernel of a sum of two commuting locally nilpotent derivations is finitely
generated when both kernels are \cite[Question~11.11, p.~291]{Freudenburg2017}.

Section~\ref{sec:commutative:jacobian} follows another open problem of
polynomial rings, the Jacobian
conjecture~\cite[p.~91]{VanEssenKurodaKoenig2021}, and asks what that
conjecture says about derivations.
% claim-ok: own signpost to the next section, moved after the problems

The LFED conjecture asks when the Lie algebra generated by a set of locally
finite derivations is itself locally finite. Lv and Yan settle it for a class
of $E$-derivations \cite[Abstract, p.~1]{LvYan2020LFED}.
